\subsection{Spectral Radius and Matrix Powers}

\begin{definition} \label{def:spectral-radius-and-convergent-matrix} \hlt{Spectral Radius and Convergent Matrix}\\
For $A\in\operatorname{Mat}(n,n,\mathbb{C})$, the \hlt{spectral radius} is $\rho(A)\defeq\max\{\abs{\lambda}:\lambda\in\operatorname{spec}(A)\}$.\\
We say $A^p\to0$ if every entry of $A^p$ converges to $0$ as $p\to\infty$.
\end{definition}

\begin{theorem} \label{thm:matrix-powers-converge-to-zero-spectral-radius} \hlt{Matrix Powers Converge to Zero}\\
Let $A\in\operatorname{Mat}(n,n,\mathbb{C})$. Then $A^p\to0$ if and only if $\rho(A)<1$.
\begin{proof}
Choose a Jordan form $A=SJS^{-1}$.\\
Since $S$ and $S^{-1}$ are fixed, $A^p\to0$ if and only if every Jordan block power in $J^p$ converges to $0$.
\begin{enumerate}[label=\roman*.]
\setlength{\itemsep}{0pt}
\item[$(\Rightarrow)$] Each block power tends to $0$, so Proposition \ref{prop:powers-of-jordan-block} gives $\abs{\lambda}<1$ for every eigenvalue. Thus $\rho(A)<1$.
\item[$(\Leftarrow)$] If $\rho(A)<1$, then every eigenvalue satisfies $\abs{\lambda}<1$. Proposition \ref{prop:powers-of-jordan-block} makes each block power tend to $0$.
\end{enumerate}
\hlt{Structural Commentary}: Similarity reduces matrix-power convergence to Jordan block-power convergence.\\
\hlt{Key Insight}: The largest eigenvalue modulus controls the asymptotic size of all blocks.
\end{proof}
\end{theorem}

\begin{figure}[H]
  \centering
  \begin{tikzpicture}[
  x=0.86cm,
  y=0.86cm,
  font=\small,
  box/.style={draw=blue!65!black, fill=blue!7, rounded corners=2pt, minimum width=2.45cm, minimum height=0.68cm, align=center},
  good/.style={draw=green!45!black, fill=green!8, rounded corners=2pt, minimum width=2.0cm, minimum height=0.62cm, align=center},
  bad/.style={draw=red!60!black, fill=red!7, rounded corners=2pt, minimum width=2.0cm, minimum height=0.62cm, align=center},
  arrow/.style={-{Stealth[length=2mm]}, thick, gray!75}
]
  \node[box] (jp) at (0,0) {$J_k(\lambda)^p$};
  \node[box] (formula) at (3.7,0) {$\binom{p}{r}\lambda^{p-r}$ terms};
  \node[good] (inside) at (7.2,0.75) {$\abs{\lambda}<1$\\decays};
  \node[bad] (outside) at (7.2,-0.75) {$\abs{\lambda}\geq1$\\does not decay};

  \draw[arrow] (jp) -- (formula);
  \draw[arrow] (formula.east) -- (inside.west);
  \draw[arrow] (formula.east) -- (outside.west);

  \draw[gray!60] (10,0) circle (0.9);
  \fill[gray!65] (10,0) circle (1.2pt) node[below, black] {\scriptsize $0$};
  \fill[green!60!black] (9.65,0.35) circle (2pt) node[above left, black] {\scriptsize inside};
  \fill[red!65!black] (10.9,0) circle (2pt) node[right, black] {\scriptsize on $\abs{\lambda}=1$};
\end{tikzpicture}

  \caption{The gray circle is $\abs{\lambda}=1$: inside it, powers decay; on or outside it, they do not.}
  \label{fig:jordan-spectral-radius-convergence}
\end{figure}

\begin{definition} \label{def:power-bounded-matrix} \hlt{Power Bounded Matrix}\\
A matrix $A\in\operatorname{Mat}(n,n,\mathbb{C})$ is \hlt{power bounded} if for each entry position $(i,j)$, the sequence of entries of $A^p$ in that position is bounded for $p=1,2,\ldots$.
\end{definition}

\begin{theorem} \label{thm:matrix-power-bounded-jordan-criterion} \hlt{Power Bounded Jordan Criterion}\\
Let $A\in\operatorname{Mat}(n,n,\mathbb{C})$. The following are equivalent:
\begin{enumerate}[label=(\roman*)]
\setlength{\itemsep}{0pt}
\item $A$ is power bounded.
\item $\rho(A)\leq1$, and every eigenvalue $\lambda$ with $\abs{\lambda}=1$ has only $1\times1$ Jordan blocks.
\item $\rho(A)\leq1$, and $\dim E_{\lambda}(A)=m_A(\lambda)$ for every $\lambda\in\operatorname{spec}(A)$ with $\abs{\lambda}=1$.
\end{enumerate}
\begin{proof}
Work in Jordan form.
\begin{enumerate}[label=(\roman*)]
\setlength{\itemsep}{0pt}
\item[$(i)$]$\Rightarrow(ii)$: If $\abs{\lambda}>1$, the diagonal entries $\lambda^p$ are unbounded.\\
If $\abs{\lambda}=1$ and $J_k(\lambda)$ has $k\geq2$, then the first superdiagonal contains $p\lambda^{p-1}$, also unbounded.
\item[$(ii)$]$\Rightarrow(i)$: Blocks with $\abs{\lambda}<1$ converge to $0$ by Proposition \ref{prop:powers-of-jordan-block}; $1\times1$ blocks with $\abs{\lambda}=1$ are bounded. Finite direct sums and fixed similarities preserve boundedness.
\item[$(ii)$]$\Leftrightarrow(iii)$: For fixed $\lambda$, $\dim E_{\lambda}(A)$ counts Jordan blocks, while $m_A(\lambda)$ counts their total size.\\
Equality holds exactly when every block has size $1$.
\end{enumerate}
\hlt{Structural Commentary}: Unit-circle eigenvalues are harmless only when their nilpotent parts vanish.\\
\hlt{Key Insight}: Nontrivial Jordan blocks on the unit circle create polynomial growth.
\end{proof}
\end{theorem}
